\chapter{الأقفال والتزامن}
\label{CH:LOCK}

في الحاسوب متعدد الأنوية (Multi-core)، تنفذ الأنوية المختلفة أجزاءً من الكود بالتوازي.
عندما تحاول أكثر من نوات معالجة الوصول إلى هيكل بيانات مشتركة في الذاكرة وتعديلها في آنٍ واحد، تنشأ المشاكل المعقدة للتزامن وحالات التسابق.

توفر \textbf{الأقفال} (Locks) آلية الاستبعاد المتبادل (Mutual Exclusion) لحماية الهياكل والموارد المشتركة، وتضمن تنفيذ الأقسام الحرجة بواسطة نوات معالجة واحدة فقط في أي لحظة زمنية.

يقدم هذا الفصل مفاهيم حالات التسابق، أقفال التناوب (Spinlocks)، ترتيب الأقفال وحالات الجمود (Deadlocks)، وأقفال النوم (Sleep-locks).

\section{حالات التسابق}
\label{sec:races}

\textbf{حالة التسابق} (Race Condition) هي خلل تزامن يتوقف فيه الناتج النهائي للبرنامج أو حالة الذاكرة على تسلسل وترتيب الأنوية في قراءة وتعديل البيانات المشتركة.

مثال: لنفترض وجود دالة تحرير الذاكرة \texttt{kfree()} التي تضيف صفحة إلى قائمة الذاكرة الحرة \texttt{freelist}:
\begin{lstlisting}
struct run *r = (struct run*)pa;
r->next = kmem.freelist;
kmem.freelist = r;
\end{lstlisting}
إذا حاولت نواتان تنفيذ هذا الكود في آنٍ واحد بدون أقفال، فقد تقرأ كل منهما نفس القيمة لـ \texttt{kmem.freelist}. ثم تقوم النواة الأولى بنسخ القيمة وتعديل الرأس، تليها النواة الثانية لتمحو تعديل النواة الأولى كليًا، مما يتسبب في تسريب الذاكرة (Memory Leak) وإفساد الهيكل.

تُسمى المقطع البرمجي المحتوي على قراءة وتعديل بيانات مشتركة بـ \textbf{القسم الحرج} (Critical Section).

\section{الكود البرمجي: الأقفال}
\label{sec:code_locks}

يقدم \texttt{xv6} نوعين من الأقفال: \textbf{قفل التناوب} (Spinlock) و \textbf{قفل النوم} (Sleeplock).

تُمثَّل أقفال التناوب في \texttt{kernel/spinlock.h}:
\begin{lstlisting}
struct spinlock {
  uint locked;       // 0 إذا كان القفل فارغًا، 1 إذا كان القفل محوزًا
  char *name;        // اسم القفل للتنقيح
  struct cpu *cpu;   // وحدة المعالجة التي تحوز القفل حاليًا
};
\end{lstlisting}

لضمان الذرية، تستخدم النواة تعليمة ذرية عتادية يقدمها معالج RISC-V تُدعى \texttt{amoswap} (Atomic Memory Swap).
تنفذ الدالة \texttt{acquire()} بالاستعانة بالدالة الذرية المضمنة في C:
\begin{lstlisting}
void
acquire(struct spinlock *lk)
{
  push_off(); // تعطيل المقاطعات أولاً
  if(holding(lk))
    panic("acquire");

  while(__atomic_test_and_set(&lk->locked, __ATOMIC_ACQUIRE) != 0)
    ; // الانتظار النشط حتى يفرغ القفل

  __sync_synchronize(); // حاجز الذاكرة (Memory Barrier)
  lk->cpu = mycpu();
}
\end{lstlisting}

أما الدالة \texttt{release()} فتقوم بتحرير القفل وتصفير راية \texttt{locked} بشكل ذري وإعادة تفعيل المقاطعات عبر \texttt{pop\_off()}.

\section{الكود البرمجي: استخدام الأقفال}
\label{sec:code_using_locks}

تحمي الأقفال كافة هياكل بيانات \texttt{xv6} المشتركة:
\begin{itemize}
  \item \texttt{kmem.lock}: حماية قائمة الذاكرة الفيزيائية الحرة.
  \item \texttt{bcache.lock}: حماية ذاكرة التخزين المؤقت للكتل (Buffer Cache).
  \item \texttt{ptable.lock} أو \texttt{p->lock}: حماية أجزاء هيكل العملية وقائمة العمليات.
  \item \texttt{cons.lock}: حماية مخزن شاشة المنصة.
\end{itemize}

\section{الجمود وترتيب الأقفال}
\label{sec:deadlock_lock_ordering}

\textbf{الجمود} (Deadlock) يحدث عندما تحاول نواتان حيازة قفلين بالترتيب نفسه المعكوس.
على سبيل المثال:
\begin{itemize}
  \item تحوز النواة A القفل $L_1$ وتطلب حيازة القفل $L_2$.
  \item تحوز النواة B القفل $L_2$ وتطلب حيازة القفل $L_1$.
\end{itemize}
ستنتظر كل نوات معالجة إلى الأبد حتى تحرر الأخرى قفلها، مما يتسبب في تعليق الكود والنظام كليًا.

لتجنب حالات الجمود، ينص القوانين في تصميم الأنظمة على \textbf{تحديد ترتيب موحد وثابت للأقفال} (Lock Ordering) يجب أن تلتزم به كافة الوظائف داخل النواة.

\section{الأقفال والمقاطعات}
\label{sec:locks_and_interrupts}

إذا كانت نوات المعالجة تحوز قفل تناوب (مثل \texttt{cons.lock})، ووُقِعت مقاطعة على نفس النواة وحاول معالج المقاطعة حيازة \texttt{cons.lock}، ستدخل النواة في حلقة انتظار نشطة غير منتهية مع نفسها مما يسبب الجمود.

لتجنب هذه الحالة، تُعطل النواة المقاطعات على النواة الحالية فورًا عبر \texttt{push\_off()} قبل تنفيذ كود الانتظار في \texttt{acquire()}، وتستعيد حالة المقاطعات عبر \texttt{pop\_off()} بعد \texttt{release()}.

\section{ترتيب التعليمات والذاكرة}
\label{sec:instruction_memory_ordering}

تقوم المعالجات المتقدمة والمترجمات (Compilers) بإعادة ترتيب التعليمات القراءة والكتابة لتوسيع كفاءة التنفيذ (Instruction Reordering).
ولضمان عدم إعادة ترتيب قراءة أو كتابة البيانات الخاضعة للحماية إلى خارج القسم الحرج المحمي بـ \texttt{acquire()} و \texttt{release()}، تستخدم النواة حواجز الذاكرة (\texttt{\_\_sync\_synchronize()} أو \texttt{fence}).

\section{البرمجة بالتعليمات الذرية}
\label{sec:programming_atomics}

تتيح بعض المعماريات هياكل بيانات غير قابلة للإغلاق (Lock-Free Data Structures) يعتمد تحويرها على التعليمات الذرية مثل Compare-And-Swap (CAS). ورغم أدائها العالي، إلا أن تعقيدها وصعوبة إثبات صحتها يفسر تفضيل الأقفال الصريحة في تصميم النواة.

\section{أقفال النوم}
\label{sec:sleep_locks}

أقفال التناوب (Spinlocks) غير مناسبة للعمليات التي تحتاج إلى الاحتفاظ بالقفل لفترات زمنية طويلة (مثل العمليات التي تقرأ البيانات من القرص الصلب)؛ لأن الانتظار النشط يستهلك طاقة وحدة المعالجة لمنع العمليات الأخرى من العمل.

يوفر \texttt{xv6} \textbf{أقفال النوم} (Sleep-locks) الممثلة بـ \texttt{struct sleeplock}:
تسمح \texttt{acquiresleep()} للعملية بالدخول في حالة النوم (\texttt{sleep()}) والتخلي عن وحدة المعالجة المركزية لعملية أخرى، ثم الاستيقاظ عبر \texttt{releasesleep()} عند تحرير القفل.

\section{العالم الحقيقي}
\label{sec:real_world_lock}

تستخدم الأنظمة التشغيلية المتقدمة مثل Linux تقنيات تزامن عالية مثل أقفال القراءة والكتابة (Read-Write Locks)، وآلية RCU (Read-Copy-Update) التي تتيح للقراء قراءة هياكل البيانات دون الحاجة لإيقاف المكاتبين أو حيازة أقفال تقليدية.

\section{تمارين}
\label{sec:exercises_lock}

\begin{enumerate}
  \item ما الفرق الأساسي بين قفل التناوب \texttt{spinlock} وقفل النوم \texttt{sleeplock}؟ ومتى يُستخدم كل منهما؟
  \item اكتب برنامج اختبار يُولّد حالة تسابق بين عمليتين وأثبت القضاء عليها باستخدام الأقفال.
\end{enumerate}
